\section{Introduction}\label{sec:introduction}

Paper I~\cite{ShcherbakovI} asks what an incomplete action specification
forces: which carrier elements become equal, which missing values are
already determined, and which quotients admit carrier-valued completions.
The present paper asks when those consequences become visible if the
syntax available to a proof is bounded. This is a question about the
chosen presentation, not just its final equality relation.

A small example captures the distinction. Let
\[
E=\{p=f^2(r),\ q=f^2(r)\}.
\]
Then $p=q$ follows through a term of depth two. At horizon zero or one,
neither input equation is available. Adding the derived equation $p=q$
changes no full equational consequence, but makes that equality available
already at horizon zero. The full theory is the same; its visibility
filtration is different.

There is a second distinction that must be fixed before measuring a proof.
One can bound the terms occurring in a proof DAG that uses reflexivity,
symmetry, transitivity, and congruence. Alternatively, one can bound the
whole terms in a sequential derivation, where each step replaces one
instance of an input equation inside a context. These proof formats give
the same unrestricted equality, but their bounded versions differ.
For the example above, a congruence proof of $f(p)=f(q)$ first proves
$p=q$ at depth two and then applies $f$ to the endpoints. Its maximum term
depth is two. A sequential derivation must pass through $f^3(r)$ and has
minimum depth three. Section~\ref{sec:rewrite} makes this separation exact.

We use $\lambda_E$ and $\delta_E$ for the congruence-proof filtration,
which is the resource computed by the finite-ground algorithm and used in
Paper I. A superscript $\mathrm{rw}$ always denotes sequential rewrite
visibility. The distinction is also necessary in the universal regime:
semigroup derivational space measures intermediate words in a rewrite
sequence, and transfers to the latter invariant.

The paper has four main contributions.
\begin{enumerate}
\item \textbf{A complete finite-ground spectrum.} Classical ground locality
bounds $\delta_E(d)$ by $\max\{d,h(E)\}$. Together with monotonicity, it
leaves a finite low-depth profile. Theorem~\ref{thm:catalan} shows that all
admissible profiles occur and that their number is Catalan. A finite
extension by derived equations can eliminate the overhead altogether;
the size of that extension remains a separate issue.
\item \textbf{Exact finite-query certificates.} Theorem~\ref{thm:querysubterm}
reduces every threshold query to one fixed input/query DAG. Incremental
activation gives the exact first-equality levels. Chronological explanation
capture prevents circular proofs, while finite threshold models certify
that an equality was unavailable below its claimed optimum. Neither
chronology alone nor a positive proof alone establishes minimality.
\item \textbf{A structural application from Paper I.} On its named action
interface, the first two quotients are the sets $X_1,X_2$ described in
Paper I, Theorem~6.6. We give the proof and lower separator for the common
four-element example, whose class counts are $N_1=5$ and $N_2=3$.
This also distinguishes visibility depth from the number of quotient
feedback rounds.
\item \textbf{The universal and sequential boundaries.} For finite
universal presentations, a computable majorant for either visibility
function is equivalent to a decidable ground word problem.
For sequential action-chain derivations, the typed invariant is exactly
semigroup space. Theorem~\ref{thm:envelope} shows that arbitrary one-sorted
contexts replace it by the least majorant with nondecreasing overhead.
The semigroup realization results are stated for $\delta^{\mathrm{rw}}$.
\end{enumerate}

Ground congruence closure, incremental explanations, and semigroup
space are established tools~\cite{NelsonOppen1980,DowneySethiTarjan1980,
NieuwenhuisOliveras2005,Olshanskii2013}. Our contribution is the exact
filtration and spectrum, its two-sided certificate interpretation, and
the precise contextual envelope for sequential action derivations.
The resource convention is part of every claim. In particular, no
finite-ground algorithm is asserted to decide arbitrary universal
presentations.

The planned Paper III develops computational representations and
checkable certificates for the project. This paper supplies its general
finite-ground reference results. The attached visibility prototype is
separate from the public QKF-Certifier implementation; the latter is not
assumed to implement every invariant or universal theory considered here.
